% Part: computability
% Chapter: machines-computations
% Section: configuration

\documentclass[../../../include/open-logic-section]{subfiles}

\begin{document}

\olfileid{cmp}{tur}{con}
\olsection{格局与计算}

\begin{explain}
回想前面我们曾追踪过“偶数机”的格局。由磁带、读写头和Turing机程序组成的假想机制，其实只是一种直观呈现Turing机计算过程的方式。严格来说，我们可以将Turing机在给定输入上的计算定义为一个\emph{格局}序列——而格局本身又由一个符号序列（对应计算过程中某时刻磁带的内容）、一个指示读写头位置的数以及一个状态组成。基于这些概念，我们就可以定义Turing机~$M$ 在给定输入上计算了什么。
\end{explain}

\begin{defn}[格局]
Turing机 $M = \tuple{Q, \Sigma, q_0,
\delta}$ 的一个\emph{格局}是一个三元组 $\tuple{C, m, q}$，其中
\begin{enumerate}
\item $C \in \Sigma^*$ 是 $\Sigma$ 上的有限符号序列；
\item $m \in \Nat$ 是满足 $m < \len{C}$ 的数；
\item $q \in Q$。
\end{enumerate}
直观上，序列~$C$ 是磁带的内容（即从最左方格到最后一个非空白或曾被访问的方格上的所有符号），$m$~是读写头正在扫描的方格编号（最左方格编号为~$0$），$q$ 是机器的当前状态。
\end{defn}

\begin{explain}
Turing机的可能输入是一个符号序列，通常是以某种形式编码数字的序列。Turing机的\emph{初始格局}是指我们启动Turing机处理该输入时所处的格局：磁带上首先是带端标记，紧接在其右侧的方格上写有输入串；读写头正扫描输入串最左侧的方格（即左端标记右侧的方格）；且机器处于指定的起始状态~$q_0$。
\end{explain}

\begin{defn}[初始格局]
Turing机 $M$ 对于输入 $I \in \Sigma^*$ 的\emph{初始格局}是
\[
\tuple{\TMendtape \frown I, 1, q_0}.
\]
\end{defn}

\begin{explain}
符号~$\frown$ 表示\emph{拼接}——输入串紧接在左端标记右侧开始。
\end{explain}

\begin{defn}
我们称格局 $\tuple{C, m, q}$ \emph{在一步内（根据~$M$）生成格局 $\tuple{C', m', q'}$}，
当且仅当
\begin{enumerate}
\item $C$ 的第 $m$ 个符号是 $\sigma$，
\item $M$ 的指令集指定了 $\delta(q, \sigma) =
  \tuple{q', \sigma', D}$，
\item $C'$ 的第 $m$ 个符号是 $\sigma'$，且
\item
\begin{enumerate}
\item 若 $D = L$，且当 $m>0$ 时 $m' = m - 1$，否则 $m'=0$；或者
\item $D = R$ 且 $m' = m + 1$；或者
\item $D = N$ 且 $m' = m$，
\end{enumerate}
\item 若 $m' = \len{C}$，则 $\len{C'} = \len{C} + 1$ 且 $C'$ 的第 $m'$ 个符号是~$\TMblank$；否则 $\len{C'}=\len{C}$。
\item 对于所有满足 $i < \len{C}$ 且 $i \neq m$ 的 $i$，$C'(i) = C(i)$。
\end{enumerate}
\end{defn}

\begin{defn}\ollabel{defn:run-output}
$M$ 在输入~$I$ 上的一个\emph{运行}是 $M$ 的格局序列 $C_i$，其中 $C_0$ 是 $M$ 在输入~$I$ 上的初始格局，且每个 $C_i$ 在一步内生成 $C_{i+1}$。

我们称 $M$ \emph{在输入 $I$ 上于 $k$ 步后停机}，如果 $C_k =
\tuple{C, m, q}$，$C$ 的第 $m$ 个符号是~$\sigma$，且 $\delta(q,
\sigma)$ 无定义。此时，$M$ 对输入~$I$ 的\emph{输出}是~$O$，其中 $O$ 是一个不以~$\TMblank$ 结尾的符号串，且存在某个~$j \in \Nat$ 使得 $C = \TMendtape \concat O \concat \TMblank^j$。（$\TMblank^j$ 表示由 $j$ 个 $\TMblank$ 组成的序列。）
\end{defn}

\begin{explain}
根据此定义，$M$ 的输出~$O$ 总是以除 $\TMblank$ 以外的符号结尾，或者，如果在时刻~$k$ 整条磁带（除最左端的~$\TMendtape$ 外）都被~$\TMblank$ 填满，则 $O$ 是空串。
\end{explain}

\end{document}